\documentclass[10pt,a4paper]{article}
\usepackage[utf8]{inputenc}
\usepackage[vietnamese]{babel}
\usepackage{amsmath,amssymb,amsfonts}
\usepackage{geometry}
\geometry{a4paper, left=10mm, right=10mm, top=10mm, bottom=10mm}
\usepackage{tikz}
\usetikzlibrary{shapes,arrows.meta,positioning,calc}
\usepackage[most]{tcolorbox}
\usepackage{array,booktabs,multicol}
\usepackage{fontawesome}
\usepackage{helvet}
\renewcommand{\familydefault}{\sfdefault}
\usepackage{parskip}

\definecolor{primary}{RGB}{112, 48, 160}
\definecolor{teal}{RGB}{0, 128, 128}
\definecolor{orange}{RGB}{230, 108, 44}
\definecolor{green}{RGB}{46, 139, 87}
\definecolor{rose}{RGB}{192, 0, 0}

\pagestyle{empty}

\begin{document}

% --- HEADER BANNER ---
\begin{tcolorbox}[
  enhanced,
  colback=violet!8!white,
  colframe=violet!80!black,
  arc=2.5mm, boxrule=1pt,
  left=3mm, right=3mm, top=2mm, bottom=2mm
]
  \begin{tabular*}{\textwidth}{@{\extracolsep{\fill}} lr}
    {\small\bfseries\color{violet!85!black}\faBook\ SỔ TAY KIẾN THỨC TOÁN 2} & 
    {\small\bfseries\color{teal!85!black}\faGraduationCap\ CHƯƠNG TRÌNH GDPT 2018}
  \end{tabular*}
  \vspace{-1mm}
  \begin{center}
    \rule{0.98\textwidth}{0.6pt}
  \end{center}
  \vspace{-2mm}
  \centering
  {\LARGE\bfseries\color{violet!90!black}\faBookmark\ BÀI 64: THU THẬP, PHÂN LOẠI, KIỂM ĐẾM SỐ LIỆU}\\[1mm]
  {\itshape\footnotesize\color{black!75} Quy trình ba bước làm quen với thống kê toán học sơ cấp}
\end{tcolorbox}

\vspace{-2mm}

\noindent
\begin{minipage}[t]{0.49\textwidth}\vspace{0pt}
\begin{tcolorbox}[
  enhanced,
  colback=violet!4!white,
  colframe=violet!75!black,
  arc=2mm, boxrule=0.8pt,
  left=2mm, right=2mm, top=1.8mm, bottom=1.8mm,
  title={\footnotesize\bfseries\color{white}\faKey\ 1. Các khái niệm ban đầu}
]
{\footnotesize
\begin{itemize}\setlength{\itemsep}{1pt}\setlength{\leftskip}{-2.5mm}
  \item \textbf{Thu thập:} Quan sát, tìm kiếm, tập hợp các đối tượng, sự vật cần tìm hiểu xung quanh ta.
  \item \textbf{Phân loại:} Sắp xếp các đối tượng đã thu thập thành từng nhóm theo các tiêu chí cụ thể:
  \begin{itemize}\setlength{\itemsep}{0.5pt}\setlength{\leftskip}{-2mm}
    \item Theo hình dạng (hình tròn, vuông, tam giác...).
    \item Theo màu sắc (xanh, đỏ, vàng...).
    \item Theo loại (hoa, quả, con vật, phương tiện...).
  \end{itemize}
  \item \textbf{Kiểm đếm:} Đếm số lượng của từng nhóm đối tượng sau khi đã phân loại.
\end{itemize}
}
\end{tcolorbox}
\end{minipage}
\hfill
\begin{minipage}[t]{0.49\textwidth}\vspace{0pt}
\begin{tcolorbox}[
  enhanced,
  colback=teal!4!white,
  colframe=teal!75!black,
  arc=2mm, boxrule=0.8pt,
  left=2mm, right=2mm, top=1.8mm, bottom=1.8mm,
  title={\footnotesize\bfseries\color{white}\faSliders\ 2. Kĩ năng ghi chép số liệu}
]
{\footnotesize
\begin{itemize}\setlength{\itemsep}{1pt}\setlength{\leftskip}{-2.5mm}
  \item \textbf{Cách đếm tránh sót, tránh trùng:}
  \begin{itemize}\setlength{\itemsep}{0.5pt}\setlength{\leftskip}{-2mm}
    \item Đếm đến đâu dùng bút chì đánh dấu gạch chéo hoặc chấm nhẹ vào hình đó.
    \item Dùng bảng kiểm đếm: mỗi đối tượng gạch một vạch ($|$), đủ $5$ vạch tạo thành một cụm gạch chéo.
  \end{itemize}
  \item \textbf{Lập bảng số liệu:} Ghi tên nhóm và số lượng tương ứng vào từng cột/hàng.
  \item \textbf{Nhận xét số liệu:} Nhóm nào nhiều nhất, nhóm nào ít nhất, nhóm nào có số lượng bằng nhau.
\end{itemize}
}
\end{tcolorbox}
\end{minipage}

\vspace{-1.5mm}
\begin{tcolorbox}[
  enhanced,
  colback=orange!4!white,
  colframe=orange!80!black,
  arc=2mm, boxrule=0.8pt,
  left=2mm, right=2mm, top=1.5mm, bottom=1.5mm,
  title={\footnotesize\bfseries\color{white}\faSitemap\ 3. Sơ đồ quy trình 3 bước xử lí số liệu thống kê}
]
\centering
\begin{tikzpicture}[x=1.65cm, y=0.75cm]
  \tikzset{arr/.style={-{Stealth[scale=0.8]}}, every node/.append style={align=center}}
  \node[rectangle, rounded corners=2pt, draw=violet!80!black, fill=violet!15, font=\scriptsize\bfseries] (b1) at (0,0) {BƯỚC 1: THU THẬP\\Quan sát đồ vật};
  \node[rectangle, rounded corners=2pt, draw=teal!80!black, fill=teal!15, font=\scriptsize\bfseries] (b2) at (2.4,0) {BƯỚC 2: PHÂN LOẠI\\Theo tiêu chí};
  \node[rectangle, rounded corners=2pt, draw=orange!80!black, fill=orange!15, font=\scriptsize\bfseries] (b3) at (4.8,0) {BƯỚC 3: KIỂM ĐẾM\\Ghi chép số lượng};

  \draw[arr, thick] (b1) -- (b2);
  \draw[arr, thick] (b2) -- (b3);
\end{tikzpicture}
\end{tcolorbox}

\vspace{-1.5mm}
\noindent
\begin{minipage}[t]{0.49\textwidth}\vspace{0pt}
\begin{tcolorbox}[
  enhanced,
  colback=green!4!white,
  colframe=green!60!black,
  arc=2mm, boxrule=0.8pt,
  left=2mm, right=2mm, top=1.8mm, bottom=1.8mm,
  title={\footnotesize\bfseries\color{white}\faCheckCircle\ 4. Bảng kiểm đếm minh họa}
]
{\footnotesize
Kiểm đếm các loại hoa trong vườn:\\
\begin{tabular}{|l|l|c|}
  \hline
  \textbf{Loại hoa} & \textbf{Kiểm đếm} & \textbf{Số lượng} \\
  \hline
  Hoa hồng & $||||$ & $\mathbf{4}$ \\
  Hoa cúc & $||||| \ ||$ & $\mathbf{7}$ \\
  Hoa sen & $|||||$ & $\mathbf{5}$ \\
  \hline
\end{tabular}
\vspace{1mm}\\
- Hoa có nhiều nhất là: \textbf{Hoa cúc} ($7$ bông).\\
- Hoa có ít nhất là: \textbf{Hoa hồng} ($4$ bông).\\
- Tổng cộng có: $4 + 7 + 5 = \mathbf{16}$ bông hoa.
}
\end{tcolorbox}
\end{minipage}
\hfill
\begin{minipage}[t]{0.49\textwidth}\vspace{0pt}
\begin{tcolorbox}[
  enhanced,
  colback=rose!4!white,
  colframe=rose!70!black,
  arc=2mm, boxrule=0.8pt,
  left=2mm, right=2mm, top=1.8mm, bottom=1.8mm,
  title={\footnotesize\bfseries\color{white}\faExclamationTriangle\ 5. Lỗi thường gặp khi đếm}
]
{\footnotesize
\begin{itemize}\setlength{\itemsep}{1.5pt}\setlength{\leftskip}{-2.5mm}
  \item \textbf{Đếm trùng lặp:} Một vật đếm đến hai lần do không đánh dấu.
  \item \textbf{Bỏ sót đối tượng:} Nhìn lướt qua mắt bỏ sót những vật nằm ở góc hoặc bị che khuất.
  \item \textbf{Phân loại sai tiêu chí:} Để lẫn đồ vật của nhóm này sang nhóm khác.
  \item Luôn cộng tổng số lượng các nhóm xem có khớp với tổng số đối tượng ban đầu không.
\end{itemize}
}
\end{tcolorbox}
\end{minipage}


\end{document}
